% part: intuitionistic-logic
% chapter: propositions-as-types
% section: proof-terms
%READERNOTE{SOL6-A323: खुणित संदर्भ सुसंगत संच आहे आणि पदांची रचना विगमनाने आहे, हे आधीच्या नियमांशी जुळवले.}

\documentclass[../../../include/open-logic-section]{subfiles}

\begin{document}

\olfileid{int}{pty}{ter}

\olsection{सिद्धता-पदे}

\Log{N1i} मध्ये गृहीतकांना निरसन-खुणा म्हणून चलांची
चिन्हे देतो (उदा., $!A^x$). \Log{N2i} मध्ये प्रत्येक
$\Gamma \Sequent !A$ क्रमवर्तीच्या डावीकडील संदर्भ~$\Gamma$
हा सुसंगत चल–!!{formula} जोड्यांचा संच असतो. त्यामुळे
\Log{N2i} मधील !!a{proof} च्या प्रत्येक क्रमवर्तीवर
आतापर्यंतच्या !!{proof} ने काय !!{prove} केले आहे
(म्हणजे~$!A$) एवढीच नव्हे, तर त्या टप्प्यावर कोणती
गृहीतके खुली आहेत (म्हणजे~$\Gamma$) ही माहितीही असते.
$!A$ \emph{कसे} !!{prove} झाले याची माहितीही प्रत्येक
क्रमवर्तीत घातली तर? त्यासाठी \emph{पद-चिन्हांकित}
\Log{tN2i} पद्धती परिभाषित करता येते; तिच्या
क्रमवर्तींचे रूप $\Gamma \Sequent N : !A$ असते.
$\Gamma$ आणि~$!A$ यांचे अर्थ आधीप्रमाणेच आहेत:
$x: !B$ या रूपातील जोड्यांचा संच, आणि क्रमवर्तीवर
संपणाऱ्या उप-!!{proof} ने~$\Gamma$ मधील गृहीतकांवरून
निष्पन्न होते असे दाखवलेले !!{formula}. ~$N$ हे त्या
क्रमवर्तीवर संपणाऱ्या !!{proof} चे सांकेतिकीकरण करणारे
विशिष्ट पद असेल.

प्रत्येक क्रमवर्तीला देण्यात येणारी पदे $x$, $y$,
\dots या चलांपासून आणि आतापर्यंत झालेल्या अनुमानांचे
सांकेतिकीकरण करणाऱ्या फलन-चिन्हांपासून बनतात.
प्रत्येक अनुमान-नियमासाठी स्वतंत्र फलन-चिन्ह लागते.
!!{introduction} नियमांना अनुरूप फलन-चिन्हांना
``रचनाकार'' आणि !!{elimination} नियमांना अनुरूप चिन्हांना
``विघटक'' म्हणतात. नियमाची जितकी आधारविधाने, तितके
प्रत्येक फलन-चिन्हाचे आर्ग्युमेंट असतात. उदाहरणार्थ,
\Intro{\land} आणि \Elim{\land}[1] साठी अनुक्रमे
$c^\land$ (दोन-स्थानी) आणि $d^\land_1$ (एक-स्थानी)
ही फलन-चिन्हे वापरता येतील. \Log{tN2i} मध्ये हे
नियम असे होतील:
\[
\Axiom$\Gamma_1 \fCenter N: !A$
\Axiom$\Gamma_2 \fCenter M: !B$
\RightLabel{\Intro{\land}}
\BinaryInf$\Gamma_1, \Gamma_2 \fCenter c^\land(N, M) : !A \land !B $
\DisplayProof
\quad
\Axiom$\Gamma \fCenter N: !A \land !B$
\RightLabel{\Elim{\land}[1]}
\UnaryInf$\Gamma \fCenter d^\land_1(N) : !A$
\DisplayProof
\]
अनुमान-नियम एखादे गृहीतक निरसित करत असेल,
म्हणजे आधारविधानात असलेले खुणित !!{formula}~$x:!A$
निष्कर्षाच्या संदर्भात नसेल, तर ते सिद्धता-पदातही
दर्शवावे लागते. अशा नियमांशी संबंधित फलन-चिन्हे
त्या चलांना \emph{बद्ध} करतात. उदाहरणार्थ,
\Intro{\lif} नियम $x:!A, \Gamma \Sequent !B$ या
आधारविधानापासून $\Gamma \Sequent !A \lif !B$
या निष्कर्षापर्यंत जातो. रचनाकार~$c^{\lif}$ एक
आर्ग्युमेंट घेतो, पण चल~$x$ बद्ध होतो हेही त्याने
दर्शवले पाहिजे. खूण~$x$ असलेले गृहीतक निरसित
झाले आहे हे सिद्धता-पद अशा प्रकारे व्यक्त करते.
उदाहरणार्थ, आधारविधानाचे सिद्धता-पद~$N$ असेल,
तर निष्कर्षाचे सिद्धता-पद $c^{\lif}([x]N)$ आणि
नियम असा लिहिता येईल:
\[
\Axiom$x:!A, \Gamma \fCenter N: !B$
\RightLabel{\Intro{\lif}}
\UnaryInf$\Gamma \fCenter c^{\lif}([x]N) : !A \lif !B$
\DisplayProof
\]

!!a{proof} पूर्णपणे पुन्हा उभारण्यासाठी आणखी
काही माहिती लागते. नियमाच्या आधारविधानांतील
!!{formula}s वरून निष्कर्षातील !!{formula} पूर्णपणे
ठरत नसेल, तर ती माहिती पदात जोडली पाहिजे.
\Intro{\lif} मध्ये असे होते; म्हणून आपण
$c^{\lif}([x:!A]N)$ असे लिहू. \Intro{\lor}
आणि~\FalseInt मध्येही असेच होते; म्हणून ती
माहिती वाहून नेणारी फलन-चिन्हे वापरू, उदाहरणार्थ:
\[
\Axiom$\Gamma \fCenter N:!A$
\RightLabel{\Intro{\lor}[1]}
\UnaryInf$\Gamma \fCenter c^{\lor{!B}}_1(N):!A \lor !B$
\DisplayProof
\quad
\Axiom$\Gamma \fCenter N:\lfalse$
\RightLabel{\FalseInt}
\UnaryInf$\Gamma \fCenter d^\lfalse_{!A}(N):!A$
\DisplayProof
\]

या कल्पनांवरून क्रमवर्ती-शैलीतील नैसर्गिक निगमनाची
अशी पद्धती बांधता येते की प्रत्येक क्रमवर्तीचे रूप
$\Gamma \Sequent N : !A$ असते आणि~$N$ हे
\emph{सिद्धता-पद} असते.

अशी सिद्धता-पदे आणि तथाकथित \emph{प्रकारयुक्त लॅम्डा कलन}
यातील पदे यांचा जवळचा संबंध आहे. तो लक्षात घेऊन वरची
सर्वसाधारण रचनाकार-विघटक चिन्हे न वापरता लॅम्डा
कलनाच्या साहित्यातील चिन्हे वापरू. उदाहरणार्थ,
$c^{\lif}([x:!A]N)$ ऐवजी $\lambd[\typeof{x}{!A}][N]$,
$d^{\lif}(N,M)$ ऐवजी $(NM)$,
$c^\land(N,M)$ ऐवजी $\pair{N}{M}$, आणि
$d^\land_i(N)$ ऐवजी~$\proj{i}{N}$ लिहू.

\begin{defn}
  सिद्धता-पदे पुढील नियमांनी विगमनाने तयार होतात:
  \begin{enumerate}
  \item प्रत्येक चर~$x$ हे सिद्धता-पद आहे (स्वयंसिद्धके).
  \item $x$ हे चर, $!A$ हे !!a{formula}, आणि~$N$ हे
    सिद्धता-पद असेल, तर $\lambd[\typeof{x}{!A}][N]$
    हेही सिद्धता-पद आहे~(\Intro\lif).
  \item $N$ आणि~$M$ ही सिद्धता-पदे असतील, तर $(NM)$
    हेही सिद्धता-पद आहे~(\Elim\lif).
  \item $N$ आणि~$M$ ही सिद्धता-पदे असतील, तर
    $\pair{N}{M}$ हे सिद्धता-पद आहे~(\Intro\land).
  \item $N$ हे सिद्धता-पद असेल, तर $\proj{i}{N}$
    हेही सिद्धता-पद आहे; येथे~$i$ हे~$1$ किंवा~$2$
    असते~(\Elim\land[i]).
  \item $N$ हे सिद्धता-पद आणि~$!A$ हे सूत्र असेल,
    तर $\inj[!A]{i}{N}$ हे सिद्धता-पद आहे; येथे~$i$
    हे~$1$ किंवा~$2$ असते~(\Intro\lor[i]).
  \item $M$, $N_1$, $N_2$ ही सिद्धता-पदे आणि~$x_1$,
    $x_2$ ही चले असतील, तर
    $\dcase{M}{x_1}{N_1}{x_2}{N_2}$ हे सिद्धता-पद
    आहे~(\Elim\lor).
  \item $N$ हे सिद्धता-पद आणि~$!A$ हे !!a{formula}
    असेल, तर $\abort{!A}{N}$ हे सिद्धता-पद आहे~(\FalseInt).
  \end{enumerate}
\end{defn}

प्रत्येक सिद्धता-पद हे एखाद्या !!a{proof} चे
सिद्धता-पद असेलच असे नाही. उदाहरणार्थ,
$\pair{N}{M}$ हे पद~\Intro{\land} अनुमानावर
संपणाऱ्या सिद्धतेचे सांकेतिकीकरण करते; म्हणून तिचे
अंतिम !!{formula} $!A \land !B$ हे संयोगसूत्र असते.
हे~\Elim{\lif} नियमाचे मुख्य आधारविधान होऊ शकत
नाही; त्यामुळे $(\pair{N}{M}P)$ हे योग्य सिद्धतेचे
सिद्धता-पद होऊ शकत नाही. \Log{tN2ip} च्या नियमांचे
पालन करणारी सिद्धता-पदेच \emph{योग्य} सिद्धता-पदे
आहेत. हे नियम \olref{tab:tN2ip} मध्ये दिले आहेत.

\subfile{rules-tN2}

\end{document}
